\documentclass[11pt,a4paper]{article}

\usepackage[margin=2.5cm]{geometry}
\usepackage{booktabs}
\usepackage{xcolor}
\usepackage{colortbl}
\usepackage{tabularx}
\usepackage{hyperref}
\usepackage{enumitem}
\usepackage{titlesec}
\usepackage{fancyhdr}
\usepackage{graphicx}
\usepackage{array}
\usepackage{microtype}
\usepackage{parskip}

% ── colours ──────────────────────────────────────────────────────────────────
\definecolor{mainblue}{HTML}{2E5FA3}
\definecolor{lightblue}{HTML}{D6E4F7}
\definecolor{lightgray}{HTML}{F5F5F5}
\definecolor{darkgray}{HTML}{333333}
\definecolor{green}{HTML}{27AE60}
\definecolor{orange}{HTML}{E67E22}

% ── section formatting ────────────────────────────────────────────────────────
\titleformat{\section}
  {\large\bfseries\color{mainblue}}
  {\thesection.}{0.5em}{}
  [\vspace{-4pt}\textcolor{mainblue}{\rule{\linewidth}{0.8pt}}]

\titleformat{\subsection}
  {\normalsize\bfseries\color{darkgray}}
  {\thesubsection.}{0.5em}{}

\titlespacing{\section}{0pt}{16pt}{6pt}
\titlespacing{\subsection}{0pt}{10pt}{4pt}

% ── header / footer ───────────────────────────────────────────────────────────
\pagestyle{fancy}
\fancyhf{}
\fancyhead[L]{\small\textcolor{mainblue}{\textbf{TCR--Epitope Binding Prediction}}}
\fancyhead[R]{\small\textcolor{darkgray}{Results Summary}}
\fancyfoot[C]{\small\textcolor{darkgray}{\thepage}}
\renewcommand{\headrulewidth}{0.4pt}

% ── hyperref ──────────────────────────────────────────────────────────────────
\hypersetup{
  colorlinks=true, linkcolor=mainblue, urlcolor=mainblue,
  pdftitle={TCR-Epitope Results Summary},
  pdfauthor={Helena Martinez Rio},
}

% ─────────────────────────────────────────────────────────────────────────────
\begin{document}

% ── title block ──────────────────────────────────────────────────────────────
\begin{center}
  {\LARGE\bfseries\color{mainblue} TCR--Epitope Binding Prediction}\\[6pt]
  {\large\bfseries\color{darkgray} Baseline Results Summary}\\[8pt]
  {\small\color{gray} VT1 Project \quad|\quad Helena Mart\'inez R\'io \quad|\quad ZHAW 2025}
\end{center}

\vspace{6pt}
\textcolor{mainblue}{\rule{\linewidth}{1.5pt}}
\vspace{12pt}

% ─────────────────────────────────────────────────────────────────────────────
\section{Project Overview}

This document summarises the baseline machine learning pipeline built to predict
whether a T-cell receptor (TCR) binds a given epitope.
The goal is to establish reference metrics before moving to deep learning and
graph-based approaches.

\subsection{Data}
\begin{itemize}[leftmargin=1.4em, itemsep=2pt]
  \item \textbf{Source:} iggytop database (\texttt{biocypher/iggytop})
  \item \textbf{Size:} $\sim$300{,}000 positive TCR--epitope pairs extracted
        via \texttt{scirpy.get.airr\_context}
  \item \textbf{Sequences:} CDR3$\beta$ (VDJ$_1$ / TRB) and epitope amino-acid strings
  \item \textbf{Negatives:} random mismatched pairing (each TCR paired with
        a random epitope it did not come from)
  \item \textbf{Final dataset:} $\sim$600{,}000 samples (1:1 positive:negative ratio)
\end{itemize}

\subsection{Feature Engineering}

No hand-crafted features. Both sequences were embedded using
\textbf{ESM2} (\texttt{facebook/esm2\_t6\_8M\_UR50D}), a protein language model
pre-trained on 250M protein sequences, yielding a 320-dimensional vector per sequence.
TCR and epitope embeddings were concatenated into a \textbf{640-dimensional feature vector}.

\subsection{Models}
\begin{itemize}[leftmargin=1.4em, itemsep=2pt]
  \item \textbf{Logistic Regression} --- linear baseline
        (scikit-learn, saga solver, balanced class weights)
  \item \textbf{Random Forest} --- non-linear baseline
        (100 trees, max depth 15, balanced class weights)
\end{itemize}

% ─────────────────────────────────────────────────────────────────────────────
\section{Model Performance}

\vspace{4pt}
\begin{table}[h!]
\centering
\renewcommand{\arraystretch}{1.3}
\begin{tabularx}{\linewidth}{>{\bfseries}l X X}
\toprule
\rowcolor{mainblue}
\textcolor{white}{Metric} &
\textcolor{white}{Logistic Regression} &
\textcolor{white}{Random Forest} \\
\midrule
\rowcolor{lightblue}
AUC-ROC     & 0.559 & \textbf{0.738} \\
AUPR        & 0.548 & \textbf{0.731} \\
\rowcolor{lightblue}
F1-score    & 0.490 & \textbf{0.743} \\
Precision   & 0.498 & \textbf{0.722} \\
\rowcolor{lightblue}
Recall      & 0.481 & \textbf{0.683} \\
Accuracy    & 0.498 & \textbf{0.714} \\
\rowcolor{lightblue}
Specificity & 0.515 & \textbf{0.745} \\
\bottomrule
\end{tabularx}
\caption{Test-set performance of both baseline models. Bold = better model per metric.}
\end{table}

\subsection{Interpretation}

\begin{itemize}[leftmargin=1.4em, itemsep=3pt]
  \item Random Forest substantially outperforms Logistic Regression (AUC 0.738 vs 0.559),
        confirming that the relationship between ESM2 features and binding is \textbf{non-linear}.
  \item Logistic Regression near AUC 0.5 indicates it cannot find a linear decision
        boundary in the 640-dimensional embedding space.
  \item AUC 0.738 is a strong baseline. Literature reports AUC 0.70--0.85 for similar
        setups with random negative generation.
  \item The gap between validation AUC (0.810) and test AUC (0.738) indicates mild
        overfitting; increasing data volume or regularisation would help.
\end{itemize}

% ─────────────────────────────────────────────────────────────────────────────
\section{Embedding Structure Analysis}

UMAP projections of ESM2 embeddings were computed separately for TCR and epitope
sequences to understand what structure the language model captures prior to any
model training.

\subsection{Key Findings}

\begin{itemize}[leftmargin=1.4em, itemsep=3pt]
  \item \textbf{TCR-only UMAP:} binding and non-binding labels are fully mixed ---
        consistent with biology (TCR sequence alone does not determine binding;
        the epitope context is required).
  \item \textbf{TCR+Epitope UMAP:} distinct isolated clusters emerge, several
        of which are strongly enriched for a single pathogen or MHC class.
        This confirms that the \emph{joint} embedding captures binding-relevant structure.
  \item \textbf{Pathogen coloring:} isolated islands correspond almost entirely to one
        pathogen (e.g., all SARS-CoV-2 pairs in one region, all Influenza pairs in another).
  \item \textbf{MHC class:} MHC class I and class II pairs separate in embedding space,
        which is biologically expected since they present different peptide lengths and
        bind different TCR subtypes.
\end{itemize}

% ─────────────────────────────────────────────────────────────────────────────
\section{Feature Importance (SHAP)}

SHAP (SHapley Additive exPlanations) was computed on the Random Forest to identify
which of the 640 ESM2 dimensions contribute most to binding predictions.

\begin{itemize}[leftmargin=1.4em, itemsep=3pt]
  \item Both TCR and epitope embedding dimensions contribute;
        no single dimension dominates the predictions.
  \item Top SHAP features are distributed across both embedding halves,
        confirming that the model exploits both sequence modalities.
  \item High-SHAP dimensions correspond to latent biochemical properties
        (hydrophobicity, charge distribution) rather than identifiable
        sequence positions.
\end{itemize}

% ─────────────────────────────────────────────────────────────────────────────
\section{Limitations}

\begin{itemize}[leftmargin=1.4em, itemsep=3pt]
  \item \textbf{Negative quality:} random mismatching may accidentally produce
        true positives, adding label noise to the training data.
  \item \textbf{Data leakage risk:} a random train/test split was used; an
        epitope-aware split (\texttt{GroupShuffleSplit}) would give a stricter
        and more realistic evaluation of generalisation to unseen epitopes.
  \item \textbf{Model size:} ESM2 8M is the smallest ESM2 variant; larger models
        (35M, 150M, 650M parameters) would produce richer representations.
  \item \textbf{Relational signal:} classical ML cannot capture the relational
        structure between TCRs and epitopes across patients and datasets.
\end{itemize}

\vspace{12pt}
\textcolor{mainblue}{\rule{\linewidth}{0.8pt}}
\begin{center}
\small\color{gray}
Generated as part of the VT1 computational immunology project, ZHAW 2025.
\end{center}

\end{document}